\subsection{Schwächen und Grenzen}
\label{subsec:schwaechen-grenzen}

Zu unterscheiden sind technische Schwächen, bewusste Scope-Grenzen und nicht
verifizierte Bereiche. Nur die erste Gruppe bezeichnet unmittelbar einen
Nachteil der Umsetzung.
Scope-Grenzen erzeugen zwar Anpassungsbedarf, sind aber keine fehlgeschlagene
Anforderung. Fehlende Verifikation erlaubt noch kein positives oder negatives
technisches Urteil.

Die reguläre Klassifikation der Dateigröße ist eindeutig: 900 Code-LOC sind
zulässig und ergeben \emph{STRONG\_WARNING}; 901 Code-LOC ergeben
\texttt{CQ001 ERROR}. Die Projektdokumentation definiert 900 LOC als absolute
Obergrenze. Die Liste zulässiger Ausnahmen enthält technisch jedoch auch
\texttt{CQ001}; ein exakt passender Eintrag kann den
901-LOC-Fehler unterdrücken und den Gate-Status auf PASS setzen. Für die
vollständige Kette von 901 LOC über \texttt{CQ001} und Exception bis zum
Gate-Status fehlt ein integrierter Regressionstest. Die Policy ist damit als
absolute Grenze beabsichtigt und dokumentiert, am untersuchten Stand aber
unterdrückbar. Diese Abweichung darf weder durch eine pauschale Erfolgsaussage noch durch eine
Änderung des technischen Projekts verdeckt werden.

Weitere Exception-Grenzen liegen in der Validierung. Führende Schrägstriche
werden vor der Absolutheitsprüfung entfernt, sodass absolute POSIX- und
Windows-Pfade normalisiert statt wie dokumentiert verworfen werden. Eine
Exception für eine existierende, aber ausgeschlossene oder nicht gescannte
Datei bleibt gültig; optionale Symbole werden nicht gegen den Quelltext
geprüft. Ablaufdatum, unbekannte IDs, Duplikate und nicht vorhandene Dateien
werden dagegen getestet und als Fehler behandelt. Befristung und Begründung
unterstützen Governance, verhindern wiederholte Verlängerung oder zu breite
Excludes aber nicht automatisch; ein Reviewprozess bleibt erforderlich.

LOC, Scope-Größe und zyklomatische Komplexität sind Strukturindikatoren, keine
Messung von Verständlichkeit, Kohäsion oder fachlicher Richtigkeit.
Importanalysen erfassen direkte statisch auflösbare Beziehungen, nicht jede
dynamische, transitive oder semantische Architekturverletzung. Der
Frontend-Resolver ignoriert unbekannte Aliase; Routerlogik unterhalb der
Größenschwelle und dünne Tauri-Adapter bleiben teilweise oder vollständig
Reviewgegenstand. Eine Datei unter 900 LOC kann deshalb weiterhin strukturell
problematisch sein.

Am historischen Beobachtungsstand folgte die profilbezogene Anwendbarkeit
vorhandenen Dateien und nicht dem Manifest; der Reporter besaß keinen eigenen
SKIP-Status. Hinzu kamen die Abweichung bei der Ruff-Auflösung in generierten
Backendprofilen und veraltete öffentliche Dokumentationspfade: README und
Tooling-Guide nannten den Quality-Einstieg nicht vollständig, die
CI-Dokumentation beschrieb noch die ältere Jobstruktur. Der Dry-Run bestätigte
ausstehende Änderungen an Index und Backlinks. Diese Befunde begrenzten die
Verständlichkeit und Profilreife der ansonsten zentralen Governance. Im
v1.0.0-Freigabestand ist speziell die Ruff-Abweichung durch die stets
bereitgestellte \path{tools/.venv} behoben; der historische Befund und seine
damaligen CI-Ergebnisse bleiben davon unberührt.

Die Bündelung im Master erhöht die interne Komplexität: Änderungen an Core,
Generator oder Katalog erfordern Regressionstests über mehrere Profile.
Außerdem nutzen Frontend und Backend die gemeinsamen JSON-Schemata nicht zur
Laufzeit; Vertragsänderungen bleiben manuell. Die direkte CLI-Auswahl trennt
\texttt{optional} und \texttt{selectable} schwächer als der Dialog. Beide
Abweichungen erzeugen Wartungs- oder Fehlbedienungsrisiken. Hinzu kommen
commitabhängig widersprüchliche Nachweise: Auf \texttt{a4487ac} bestanden
Webprofile und lokale Tooling-Tests, während externe Tooling-, Desktopprofil-
und Windows-Prüfungen scheiterten. Auf \texttt{461fc751} bestand das
Master-Quality-Gate, backendhaltige Profil-Gates und Windows dagegen nicht.
Diese Befunde stellen frühere commitgebundene Abnahmen nicht rückwirkend
infrage, verhindern aber ein vollständig grünes Urteil für den
Governance-Stand.

Bewusst ausgeschlossen sind Backend in \texttt{web-only}, Cloud-API in
\texttt{desktop-local} und PostgreSQL ohne Backend. Auch Fachlogik,
Authentifizierung, Rollenmodelle, verpflichtende Persistenz, Cloud-Provider
und native Fachkommandos bleiben produktspezifisch. Der neue Entscheidungsrahmen
für Produktpersistenz stellt keine universelle lokale Datei-,
Embedded-Database- oder Synchronisationsimplementierung bereit. Diese bewusste
Scope-Grenze ist ein Extension Path, kein Implementierungsfehler
\parencite{kleiveist2026persistence}.

Nicht abschließend verifiziert sind Compose-Start, produktive Bereitstellung,
Release Validation und Branch Protection. Linux und macOS wurden unsigniert
paketiert, Windows scheiterte. Signing, Notarisierung, Veröffentlichung und
produktiver Cloudbetrieb bleiben produkt- oder infrastrukturspezifisch
\parencite{kleiveist2026governance-core-run,kleiveist2026governance-desktop-run,kleiveist2026acceptance}.
Die Paketergebnisse belegen daher weder signierte Releases noch produktiven
Cloudbetrieb. Bewertet wird die Template-Baseline, nicht die Produktionsreife
eines Produkts.
Die verbleibenden Risiken betreffen damit vor allem Betrieb, Plattformen und
Freigabeprozesse.
\beginnerinput{workspace/statement/700_bewertung/730}
